\clearpage\pagestyle{otMain}\setlength{\otBlocEspacement}{2.5mm}\fontsize{11}{15}\selectfont\corrigesChapitre{12}{Construire et comparer des triangles}
\begin{blocLivre}{cor-12-1}\begin{corrige}{1}{Vérifier un triangle}\demarche{Tester l’inégalité triangulaire.}
La plus grande longueur est 6 et $4+5=9>6$ : le triangle existe. On peut tracer d’abord le côté de 6 cm, puis utiliser les deux autres longueurs comme rayons de cercles.
\label{fig-cor-12-1}\figureGeo[corrige]{tri-01}
\end{corrige}\end{blocLivre}
\begin{blocLivre}{cor-12-2}\begin{corrige}{2}{Un triangle impossible}\demarche{Comparer le plus grand côté à la somme des autres.}
Non : la plus grande longueur 6 dépasse $2+3=5$. Les deux petits segments ne peuvent rejoindre les extrémités du grand côté. L’inégalité triangulaire n’est pas respectée.
\label{fig-cor-12-2}\figureGeo[corrige]{tri-02}
\end{corrige}\end{blocLivre}
\begin{blocLivre}{cor-12-3}\begin{corrige}{3}{Deux côtés et un angle}\demarche{Placer le sommet de l’angle connu.}
Tracer $AB=5$ cm. En A, construire une demi-droite formant $60^\circ$ avec $[AB)$. Y placer C à 4 cm de A, puis relier B à C. Les données sont deux côtés et leur angle compris : cas $\res{CAC}$.
\label{fig-cor-12-3}\figureGeo[corrige]{tri-03}
\end{corrige}\end{blocLivre}
\begin{blocLivre}{cor-12-4}\begin{corrige}{4}{Un côté et deux angles}\demarche{Utiliser la somme des angles du triangle.}
Tracer AB. En A et B, tracer les demi-droites formant les angles demandés dans le même demi-plan. Leur intersection est C. $\widehat{ACB}=180-50-60=\res{70^\circ}$. Les données correspondent au cas ACA.
\label{fig-cor-12-4}\figureGeo[corrige]{tri-04}
\end{corrige}\end{blocLivre}
\begin{blocLivre}{cor-12-5}\begin{corrige}{5}{Trois côtés correspondants}\demarche{Associer les sommets à partir des côtés communs.}
Les trois côtés sont respectivement égaux : cas CCC. On associe A à D, B à E et C à F ; les triangles ABC et DEF sont égaux. On en déduit notamment $\widehat{ABC}=\widehat{DEF}$.
\label{fig-cor-12-5}\figureGeo[corrige]{tri-05}
\end{corrige}\end{blocLivre}
\begin{blocLivre}{cor-12-6}\begin{corrige}{6}{Un angle compris}\demarche{Vérifier le rôle de l’angle avant d’utiliser CAC.}
Les angles connus sont compris entre les deux côtés égaux. Le critère CAC donne l’égalité de ABC et DEF, avec A correspondant à D. Les troisièmes côtés sont donc égaux : $\res{BC=EF}$.
\label{fig-cor-12-6}\figureGeo[corrige]{tri-06}
\end{corrige}\end{blocLivre}
\begin{blocLivre}{cor-12-7}\begin{corrige}{7}{Deux angles aux extrémités}\demarche{Citer le côté encadré par les deux angles.}
Chaque troisième angle vaut $180-50-60=\res{70^\circ}$. Le côté AB correspond à DE et les angles à ses extrémités sont égaux : cas ACA. Les triangles ABC et DEF sont donc égaux.
\label{fig-cor-12-7}\figureGeo[corrige]{tri-07}
\end{corrige}\end{blocLivre}
\begin{blocLivre}{cor-12-8}\begin{corrige}{8}{Des données insuffisantes}\demarche{Chercher une variation conservant les données.}
Non. L’angle compris entre ces côtés peut changer, ce qui change le troisième côté. Les deux positions représentées fournissent deux formes différentes malgré les mêmes deux longueurs. Il manque une donnée suffisante.
\label{fig-cor-12-8}\figureGeo[corrige]{tri-08}
\end{corrige}\end{blocLivre}
\begin{blocLivre}{cor-12-9}\begin{corrige}{9}{Un triangle isocèle}\demarche{Répartir la somme restante en deux angles égaux.}
Les angles à la base B et C sont égaux. Leur somme vaut $180-74=106^\circ$ ; chacun vaut $106/2=\res{53^\circ}$. On utilise une propriété du triangle isocèle, cohérente avec sa symétrie.
\label{fig-cor-12-9}\figureGeo[corrige]{tri-09}
\end{corrige}\end{blocLivre}
\begin{blocLivre}{cor-12-10}\begin{corrige}{10}{Un parallélogramme coupé}\demarche{Utiliser les côtés opposés et le côté commun.}
Dans un parallélogramme, $AB=CD$ et $BC=DA$. Le côté AC est commun aux deux triangles. Le cas CCC montre l’égalité de ABC et CDA, avec A correspondant à C, B à D et C à A.
\label{fig-cor-12-10}\figureGeo[corrige]{tri-10}
\end{corrige}\end{blocLivre}
\begin{blocLivre}{cor-12-11}\begin{corrige}{11}{Un cerf-volant}\demarche{Transférer un angle après avoir prouvé l’égalité.}
Les triangles ABC et DBC partagent BC. Avec $AB=DB$ et $AC=DC$, leurs trois côtés sont respectivement égaux. Le cas CCC donne l’égalité des triangles ; A correspond à D. Les angles demandés sont donc égaux.
\label{fig-cor-12-11}\figureGeo[corrige]{tri-11}
\end{corrige}\end{blocLivre}
\begin{blocLivre}{cor-12-12}\begin{corrige}{12}{Une hauteur et une aire}\demarche{Associer une base à sa hauteur perpendiculaire.}
$\mathcal A=AB\times CH/2=8\times3/2=\res{12\ \mathrm{cm}^2}$. La hauteur CH est perpendiculaire à la droite de la base AB. Le côté AC ne possède pas nécessairement cette propriété.
\label{fig-cor-12-12}\figureGeo[corrige]{tri-12}
\end{corrige}\end{blocLivre}
\begin{blocLivre}{cor-12-13}\begin{corrige}{13}{Un périmètre transféré}\demarche{Compléter les côtés grâce à la correspondance.}
Par correspondance, $BC=EF=5$ et $AC=DF=4$. Le périmètre est $7+5+4=\res{16\ \mathrm{cm}}$. Il faut respecter les sommets associés pour transférer les données.
\label{fig-cor-12-13}\figureGeo[corrige]{tri-13}
\end{corrige}\end{blocLivre}
\begin{blocLivre}{cor-12-14}\begin{corrige}{14}{Réparer un protocole}\demarche{Lire le sommet au milieu du nom de l’angle.}
Le sommet de $\widehat{BAC}$ est A, lettre du milieu. L’angle de $35^\circ$ doit donc être construit en A entre AB et la demi-droite portant AC. Placer cet angle en B ne respecte pas la donnée et ne construit pas nécessairement le triangle demandé.
\label{fig-cor-12-14}\figureGeo[corrige]{tri-14}
\end{corrige}\end{blocLivre}
\begin{blocLivre}{cor-12-15}\begin{corrige}{15}{Démontrer une perpendicularité}\demarche{Comparer deux triangles puis utiliser l’alignement.}
$AB=AC$, $BM=CM$ car M est le milieu, et AM est commun. Les triangles ABM et ACM sont égaux par CCC. Donc $\widehat{AMB}=\widehat{AMC}$. Ces deux angles adjacents forment un angle plat car B, M, C sont alignés : chacun vaut $90^\circ$. Ainsi $\res{(AM)\perp(BC)}$.
\label{fig-cor-12-15}\figureGeo[corrige]{tri-15}
\end{corrige}\end{blocLivre}
